\section*{الفصل \ref{ch:b3:many-electron-atoms} --- الذرات المتعددة الإلكترونات ورموز الحدود}

\begin{solution}{exo:b3:many-electron-atoms:1}
\[
\Psi = \frac{1}{\sqrt6}\begin{vmatrix}1s\alpha(1) & 1s\beta(1) & 2s\alpha(1)\\
1s\alpha(2) & 1s\beta(2) & 2s\alpha(2)\\ 1s\alpha(3) & 1s\beta(3) & 2s\alpha(3)
\end{vmatrix}.
\]
لا يقدّم $1s$ إلا \omterm{def:b3:many-electron-atoms:spin-orbital}{مدارين مغزليين}، $1s\alpha$ و$1s\beta$؛ وإلكترون ثالث
سيكرر أحدهما، فيتساوى عمودان: والمحدد معدوم.
\end{solution}

\begin{solution}{exo:b3:many-electron-atoms:2}
$\binom62 = 15$، $\binom63 = 20$، $\binom{10}{2} = 45$. $p^3$: \babelsublr{${}^4$S} (4) +
\babelsublr{${}^2$D} (10) + \babelsublr{${}^2$P} (6) = 20. $d^2$: \babelsublr{${}^3$F} (21) + \babelsublr{${}^3$P} (9) + \babelsublr{${}^1$G} (9) +
\babelsublr{${}^1$D} (5) + \babelsublr{${}^1$S} (1) = 45.
\end{solution}

\begin{solution}{exo:b3:many-electron-atoms:3}
N ($2p^3$، نصف ممتلئة): \termsym{4}{S}{3/2}. O ($2p^4$، أكثر من النصف):
\termsym{3}{P}{2}. F ($2p^5$): \termsym{2}{P}{3/2}. \ce{Ti^2+} ($3d^2$):
\termsym{3}{F}{2}. \ce{Fe^3+} ($3d^5$، نصف ممتلئة، كل السبينات متوازية، $L = 0$):
\termsym{6}{S}{5/2}.
\end{solution}

\begin{solution}{exo:b3:many-electron-atoms:4}
\babelsublr{${}^2$D}: $J = \frac52$ (6 حالات) و$\frac32$ (4)؛ $6 + 4 = 10 = 5 \times 2$.
\babelsublr{${}^3$P}: $J = 2$ (5)، 1 (3)، 0 (1)؛ $5 + 3 + 1 = 9 = 3 \times 3$.
\end{solution}

\begin{solution}{exo:b3:many-electron-atoms:5}
الفاصلان $16.4$ ($J = 1 - 0$) و$27.0$ ($2 - 1$): النسبة 1.65 بدل 2.
$A = 16.4/1 = \qty{16.4}{cm^{-1}}$ من الأول، و$27.0/2 = \qty{13.5}{cm^{-1}}$
من الثاني. فلا تصلح قيمة واحدة للثابت $A$: المستويات ممتزجة قليلًا بمستويات
حدود أخرى (\babelsublr{${}^1$D}، \babelsublr{${}^1$S})، \omterm{def:b3:many-electron-atoms:russell-saunders}{واقتران راسل--ساوندرز}
تقريب.
\end{solution}

\begin{solution}{exo:b3:many-electron-atoms:6}
$\lambda = 10^7/\tilde\nu$: \qty{589.76}{nm} (\babelsublr{D$_1$}) و\qty{589.16}{nm}
(\babelsublr{D$_2$}). الفاصل $\qty{17.20}{cm^{-1}} = \qty{2.13}{meV}$. $E(\frac32) -
E(\frac12) = \frac32A$، ومنه $A = \qty{11.47}{cm^{-1}}$.
\end{solution}

\begin{solution}{exo:b3:many-electron-atoms:7}
$3p \to 3s$: مسموح ($\Delta l = -1$). $3d \to 3s$: ممنوع ($\Delta l = -2$).
$3d \to 3p$: مسموح. $4s \to 3s$: ممنوع ($\Delta l = 0$، لا تغيّر في
الزوجية). $4s \to 3p$: مسموح.
\end{solution}

\begin{solution}{exo:b3:many-electron-atoms:8}
متوسط \babelsublr{${}^3$P}: $(5 \times 169086.76 + 3 \times 169086.84 + 169087.83)/9 =
\qty{169086.91}{cm^{-1}}$. الفجوة إلى \babelsublr{${}^1$P}: \qty{2047.98}{cm^{-1}} $=
\qty{0.254}{eV}$، ومنه $K(1s,2p) = \qty{0.127}{eV}$، أصغر بثلاث مرات من
$K(1s,2s)$. \omterm{def:b3:many-electron-atoms:exchange}{فتكامل التبادل} يقيس تداخل كثافتي
المدارين؛ والمدار $2s$ يتغلغل في منطقة $1s$ (فله كثافة
قرب النواة)، أما المدار $2p$ فينعدم عند النواة ويتداخل أقل.
\end{solution}

\begin{solution}{exo:b3:many-electron-atoms:9}
أكبر قيمة $M_L = 4$ تقتضي أن يكون الإلكترونان كلاهما في $m_l = 2$، بسبينين متعاكسين: \omterm{def:b3:many-electron-atoms:singlet-triplet}{حالة
أحادية}، \babelsublr{${}^1$G} (9 حالات). والتالية، $M_L = 3$، تأتي من $m_l = 2, 1$، مع
إمكان $M_S = 1$: \babelsublr{${}^3$F} (21). ويبقى $M_L = 2$ مع $M_S = 0$ وحدها:
\babelsublr{${}^1$D} (5). ويبقى $M_L = 1$ مع $M_S = 1$: \babelsublr{${}^3$P} (9). وتبقى حالة واحدة ذات
$M_L = M_S = 0$: \babelsublr{${}^1$S}. والمجموع 45.
\end{solution}

\begin{solution}{exo:b3:many-electron-atoms:10}
عند $\mathbf r_1 = \mathbf r_2 = \mathbf r$ تساوي الدالة الثلاثية
$\frac{1}{\sqrt2}[1s(\mathbf r)2s(\mathbf r) - 2s(\mathbf r)1s(\mathbf r)] = 0$؛ أما
الدالة الأحادية فتساوي $\sqrt2\,1s(\mathbf r)2s(\mathbf r) \ne 0$. ففي \omterm{def:b3:many-electron-atoms:singlet-triplet}{الحالة الثلاثية} لا
يوجد الإلكترونان أبدًا في النقطة نفسها، ويكونان في المتوسط أكثر
تباعدًا، فيكون تنافرهما أصغر: $E_+ - E_- = 2K > 0$.
\end{solution}

\begin{solution}{exo:b3:many-electron-atoms:11}
$d^8$ أكثر من نصف ممتلئة: الاقتران مقلوب، و$J = L + S = 4$
هو الأدنى. الفاصلان: $E(3) - E(4) = 1360.7$ و$E(2) - E(3) = 908.9$؛ والنسبة
1.50، مقابل $4/3 = 1.33$ عند لانديه. $|A| = 1360.7/4 = \qty{340}{cm^{-1}}$
($A < 0$). \omterm{def:b3:many-electron-atoms:spin-orbit}{فاقتران السبين--المدار} أكبر بكثير في هذا الأيون الأثقل منه في
الكربون (\qty{16}{cm^{-1}}).
\end{solution}

\begin{solution}{exo:b3:many-electron-atoms:12}
$\sum_J(2J+1)[J(J+1) - L(L+1) - S(S+1)] = 0$ (يمكن عدّ الحالات $(2L+1)(2S+1)$
في أي من الأساسين، و$\sum M_LM_S = 0$ على حالات الأساس غير المقترن).
من أجل \babelsublr{${}^3$P}: $1(0 - 4) + 3(2 - 4) + 5(6 - 4) = -4 - 6 + 10 = 0$. والمتوسط
المرجَّح للكربون $(0 + 3 \times 16.42 + 5 \times 43.41)/9 = \qty{29.6}{cm^{-1}}$:
وهي الطاقة التي كانت للحد لولا \omterm{def:b3:many-electron-atoms:spin-orbit}{اقتران السبين--المدار}.
\end{solution}

\begin{solution}{pb:b3:many-electron-atoms:1}
\textbf{1.} $1s\alpha$، $1s\beta$، $2s\alpha$، $2s\beta$.
\textbf{2.} $|1s\alpha\,2s\alpha|$، $|1s\alpha\,2s\beta|$، $|1s\beta\,2s\alpha|$،
$|1s\beta\,2s\beta|$.
\textbf{3.} $|1s\alpha\,2s\alpha| = \frac{1}{\sqrt2}[1s(1)2s(2) - 2s(1)1s(2)]\,
\alpha(1)\alpha(2)$، بنشر المحدد $2\times2$.
\textbf{4.} مجموعهما $\frac{1}{\sqrt2}[1s(1)2s(2) - 2s(1)1s(2)]\cdot
\frac{1}{\sqrt2}[\alpha(1)\beta(2) + \beta(1)\alpha(2)]$ (مضروبًا في $\sqrt2$، ثم
منظَّمًا)؛ وفرقهما يعطي $\frac{1}{\sqrt2}[1s(1)2s(2) +
2s(1)1s(2)]\cdot\frac{1}{\sqrt2}[\alpha(1)\beta(2) - \beta(1)\alpha(2)]$.
\textbf{5.} الجزء الفضائي المتناظر مع دالة السبين ضد المتناظرة ($S = 0$)؛
والجزء الفضائي ضد المتناظر مع دوال السبين المتناظرة الثلاث ($S = 1$).
\textbf{6.} $2S + 1 = 1$ و3: حالة واحدة وثلاث حالات للطاقة نفسها.
\textbf{7.} $[\ce{Ar}]\,3d^2$؛ $\binom{10}{2} = 45$ محددًا.
\textbf{8.} أكبر $S = 1$، وأكبر $L = 2 + 1 = 3$، أقل من نصف ممتلئة:
\termsym{3}{F}{2}، وهو المستوى عند 0 في المعطيات.
\textbf{9.} $21 + 9 + 9 + 5 + 1 = 45$.
\textbf{10.} عادية ($J = 2$ الأدنى). الفاصلان 184.9 و235.5: النسبة 1.27
مقابل $4/3 = 1.33$؛ فالقاعدة صحيحة في حدود 5\,\%.
\textbf{11.} \babelsublr{${}^3$F}: $(5 \times 0 + 7 \times 184.9 + 9 \times 420.4)/21 =
\qty{241.8}{cm^{-1}}$؛ \babelsublr{${}^3$P}: $(10538.4 + 3 \times 10603.6 + 5 \times
10721.2)/9 = \qty{10661.7}{cm^{-1}}$.
\textbf{12.} $15B = \qty{10419.9}{cm^{-1}}$، $B = \qty{695}{cm^{-1}}$.
\textbf{13.} $\Delta S = 0$.
\textbf{14.} $\Delta S = 1$ ممنوع، و$1s2s \to 1s^2$ له $\Delta l = 0$
(لا تغيّر في الزوجية؛ و\babelsublr{${}^3$S$_1 \to {}^1$S$_0$} له أيضًا $\Delta L = 0$ بين حالتين
S). فالحالة شبه مستقرة: تعيش مدة أطول بكثير من الحالات المثارة
العادية.
\textbf{15.} $169086.91 - 159855.97 = \qty{9230.94}{cm^{-1}}$: \qty{1083.3}{nm}، في
الأشعة تحت الحمراء القريبة.
\textbf{16.} $171134.89 - 166277.44 = \qty{4857.45}{cm^{-1}}$: \qty{2058.7}{nm}.
\textbf{17.} $\qty{171134.89}{cm^{-1}}$: \qty{58.43}{nm}، في الأشعة فوق البنفسجية
الفراغية (يمتصها الهواء).
\textbf{18.} لكل عائلة خطوطها وحالتها الدنيا، 
والعائلتان لا تتصلان بالضوء أبدًا: فهما تسلكان سلوك غازين لهما طيفان.
\textbf{19.} $E_\pm = E_{1s} + E_{2s} + J \pm K$ (الأحادية $+$، والثلاثية $-$).
\textbf{20.} الثلاثية: تنعدم دالتها الفضائية حين يلتقي الإلكترونان،
فيتنافران أقل.
\textbf{21.} \qty{6421.47}{cm^{-1}}، $6421.47/8065.54 = \qty{0.796}{eV}$.
\textbf{22.} الفجوة $\qty{2047.98}{cm^{-1}} = \qty{0.254}{eV}$، ومنه $K(1s,2p) =
\qty{0.127}{eV}$.
\textbf{23.} المدار $2s$ يتغلغل في منطقة $1s$ ويتداخل معها أكثر مما يفعل
$2p$.
\textbf{24.} نعم: في التوزيعين كليهما تقع \omterm{def:b3:many-electron-atoms:singlet-triplet}{الحالة الثلاثية} (ذات $S$ الأكبر) في الأدنى.
\textbf{25.} $K = \frac12 \times \qty{0.796}{eV}$: \textbf{\omterm{def:b3:many-electron-atoms:exchange}{تكامل التبادل}
$K(1s,2s)$ للهيليوم يساوي \qty{0.398}{eV}}، مقيسًا بوصفه نصف
الفجوة بين الحالتين الأحادية والثلاثية.
\end{solution}
